\problemname{Group Division}
Math teacher Maria is going to divide her students into two groups during the next lesson.
She is therefore now faced with a common problem: how do you divide students into two groups in a sensible way?
There are $n$ students in the class, and the classroom has $n$ chairs, numbered from $1$ to $n$.
When a student arrives in the classroom, they always sit on the leftmost empty chair.
So if a total of $k$ students show up on a certain day, they will always sit on chairs $1,2,...,k$.

To divide the students into two groups, Maria has a certain strategy she wants to use.
Before the lesson begins, she chooses a sequence $s$ of $n$ ones and twos.
Once the lesson has started, she goes to each student and assigns the student sitting on chair $i$ to group $s_i$.
For the group division to be sensible, there are two requirements that $s$ must satisfy:

\begin{enumerate}
  \item Maria does not know how many students will come, but regardless of how many show up, the difference in size between the two groups must be at most $1$.

  \item There are $m$ pairs of chairs that have the same color. If students are sitting on such a pair of chairs, they must end up in the same group.
\end{enumerate}

Your task is, given $n,m$ and the $m$ pairs of chairs, to find the sequence $s$ that comes first in alphabetical order and satisfies the requirements above.
If no valid $s$ exists, your program should print $-1$.

By alphabetical order we mean that two sequences are compared by primarily looking at the first element; if they are the same, the second element is compared, and so on.
For example, the sequence \texttt{1122} comes before \texttt{1211}, but after \texttt{1112}.

\section*{Input}
The first line contains two integers $n$ and $m$ ($1\leq n \leq 10^5$, $0 \leq m \leq \frac{n}{2}$), the number of chairs in the classroom and the number of pairs of chairs with the same color.
Then follow $m$ lines each containing two integers $s_{1}, s_{2}$ ($1 \leq s_{1}, s_{2} \leq N$, $s_{1} \neq s_{2}$), meaning that chairs $s_{1}$ and $s_{2}$ have the same color.
Each chair has the same color as at most one other chair.

\section*{Output}
Print a sequence of $n$ ones or twos (without spaces), or $-1$ if no valid solution exists.

\section*{Scoring}
Your solution will be tested on a set of test groups.
To earn points for a group, you must pass all test cases in that group.

\noindent
\begin{tabular}{| l | l | l |}
  \hline
  \textbf{Group} & \textbf{Points} & \textbf{Constraints} \\ \hline
  $1$     & $12$         & $n \leq 15$\\ \hline
  $2$     & $23$         & $n \leq 40$ \\ \hline
  $3$     & $25$         & Chair $i$ can only have the same color as chair $i\pm 1$. \\ \hline
  $4$     & $40$         & No additional constraints. \\ \hline
\end{tabular}

\section*{Explanation of Sample 1}
If only two people come, they must be in different groups.
Thus chairs $1$ and $2$ must be in different groups, but we also know that chair $3$
must be in the same group as chair $2$. If four people come, then chair $4$
must be in the same group as chair $1$ for the groups to be equal in size. But chair $4$ must
be in the same group as chair $5$. If we continue in this way, we get that chairs $1, 4$
and $5$ must be in the same group, while chairs $2, 3, 6$ and $7$ are in the other group.
There are thus two solutions: $1221122$ and $2112211$. The alphabetically smallest of
these is $1221122$.
